\clearpage
\fancyhead[L]{\small PHILOSOPHICAL TRANSACTIONS.}
\newcommand{\ObservedWollastonDenominator}{\hbox{\vrule width0pt height4.600000bp depth0pt\special{pdf:bcontent}\special{pdf:code q 0 g 2.453333 4.600000 m 2.146667 4.293333 l 2.453333 3.986667 l 2.453333 3.373333 l 1.840000 2.760000 l 1.226667 2.760000 l 1.226667 2.453333 l 0.613333 1.840000 l 0.306667 1.840000 l 0.000000 1.533333 l 0.000000 0.613333 l 0.613333 0.000000 l 3.373333 0.000000 l 3.986667 0.613333 l 3.986667 0.920000 l 4.293333 1.226667 l 4.293333 1.533333 l 3.986667 1.840000 l 3.986667 2.146667 l 3.680000 2.146667 l 3.373333 2.453333 l 3.373333 3.066667 l 3.680000 3.373333 l 3.680000 3.986667 l 3.066667 4.600000 l h 1.226667 2.146667 m 1.533333 2.453333 l 1.840000 2.453333 l 2.453333 1.840000 l 2.453333 1.226667 l 1.840000 0.613333 l 1.226667 1.226667 l h f* Q}\special{pdf:econtent}\kern4.293333bp}}

\hypertarget{native-566-lower}{}
\begin{center}\small Göttingische gelehrte Anzeigen. Stück 27. Seite 269..272. 1814 Februar 14.\\[3pt]\rule{.72\linewidth}{.3pt}\end{center}

\emph{Philosophical Transactions of the Royal Society of London} for the Jear 1813. 
IV und 304 Seiten. Meteorological Journal und 8 Seiten Index in Quart.

\emph{Beobachtung der Sommersolstitien von} 1812 \emph{und} 1813, \emph{und des Wintersolstitiums} 1812, \emph{auf der königl. Sternwarte zu Greenwich}, von J. \textsc{Pond} (drei Aufsätze). Das Instrument, womit diese Beobachtungen angestellt wurden, ist ein von \textsc{Troughton} verfertigter sechsfussiger Muralkreis, von welchem Herr \textsc{Pond} in der Folge eine ausführliche Beschreibung zu geben verspricht. Aus den wenigen Notizen, die hier davon mitgetheilt werden, sieht man, dass das Charakteristische desselben darin besteht, die Bögen des himmlischen Meridians ohne eine Beziehung auf das Zenith zu messen, daher er auch weder ein Bleiloth noch ein Niveau hat. Die Bestimmung der Stellung des Zenithpunkts im Meridian wird lediglich dem Zenithsector vorbehalten. Was hier von der Genauigkeit der Messungen mit diesem Instrument (welches einer über Paris erhaltenen Privatnachricht zufolge 1600 Guineas kosten soll) gerühmt wird, scheint eine neue Epoche für diesen Theil der practischen Astronomie zu versprechen. Die Resultate dieser Bestimmungen für die mittlere Schiefe der Ekliptik sind:

\begin{center}\fontsize{9}{11}\selectfont\setlength{\tabcolsep}{3pt}\renewcommand{\arraystretch}{1.18}\begin{tabular}{@{}l c r@{\;}r@{\;}r@{}}
Sommersolstitium von 1812 & $.\;.\;.\;.\;.\;.\;.\;.\;.\;.$ & $23^\circ$ & $27'$ & $50''5$\\
Wintersolstitium von 1812 & $.\;.\;.\;.\;.\;.\;.\;.\;.\;.$ & $23$ & $27$ & $47.35$\\
Sommersolstitium von 1813 & $.\;.\;.\;.\;.\;.\;.\;.\;.\;.$ & $23$ & $27$ & $49.5$\\
\end{tabular}\end{center}

{\rightskip=0pt plus 3em Bei Berechnung der Beobachtungen hat sich Hr. \textsc{Pond} der \textsc{Bradley}'schen Refractionstafeln bedient.\par}

\emph{Verzeichniss der Nordpolabstände von} 44 \emph{der vornehmsten Fixsterne}, für den Anfang des J. 1813, und in einem spätern Aufsatze \emph{ein ähnliches Verzeichniss für} 84 \emph{Fixsterne} von J. \textsc{Pond}. Das erstere die\SourceLeaf{567}ser Verzeichnisse ist auch in Deutschland schon früher bekannt geworden: das zweite ist als eine vermehrte und verbesserte Ausgabe zu betrachten. Die Anzahl der Beobachtungen, worauf sich diese Bestimmungen gründen, geht bei einigen Sternen bis auf 100, beim Polarstern sogar auf 200. Bei einigen der vornehmsten Sterne gibt Herr \textsc{Pond} zugleich die Resultate aus den einzelnen Beobachtungsreihen (jede von zehn Beobachtungen), deren Uebereinstimmung von der Vortrefflichkeit des Instruments zeugt und alles übertrifft, was man bisher in dieser Art gekannt hat. Bei vielen Sternen hielt Herr \textsc{Pond} seine Bestimmung auf $\frac{1}{4}$ Secunde zuverlässig, so weit sie blos vom Instrument abhängt. Eine Auswahl von 30 der vornehmsten Sterne nordwärts vom Aequator hat Herr \textsc{Pond} noch in einem besondern Verzeichniss zusammengestellt, welches er den Normalcatalog (Standard catalogue) nennt, und dem er als Seitenstück des \textsc{Maskelyne}schen Catalogs für die geraden Aufsteigungen eine immer grössere Vollkommenheit zu geben gedenkt.

Noch stehen mit der practischen Astronomie in Verbindung zwei Aufsätze von \textsc{William Hyde Wollaston}. Der erstere gibt eine Methode, äusserst feine Metalldrähte zu ziehen. Sie besteht darin, in das Innere eines anfangs dicken Silberdrahts einen Draht von Gold oder Platina zu bringen (indem man entweder den ersten hohl bohrt, oder den Platindraht in einer Form mit geschmolzenem Silber umgiesst), diesen Doppeldraht auf bekannte Weise auszudehnen, und endlich das Silber durch warme Salpetersäure wieder abzulösen. Mit Platina gelangen diese Versuche am besten. Er erhielt auf diese Weise Fäden von $\frac{1}{1000}$ und $\frac{1}{8000}$ Zoll Dicke, die zur Einspannung im Brennpunkte astronomischer Fernrohre vortreffliche Dienste zu leisten geschickt waren. Diese Fäden sind schon weit dünner als die feinsten Spinnfäden. (Rec. fand durch Messungen, dass die Dicke eines in einem 16zolligen Fernrohre im Brennpunkte eingezogene Spinnefadens einen Bogen von $7''8$ deckte; demzufolge wäre die Dicke dieses Fadens $\frac{1}{1630}$ Zoll). Allein Herr \textsc{Wollaston} trieb die Sache noch viel weiter und stellte Fäden von fast unglaublicher Feinheit (obwohl nur in kleinen Stücken) dar. Verhältnissmässig hatten sie noch viel Stärke; ein Platinafaden von $\frac{1}{18000}$ Zoll Dicke trug noch $1\frac{1}{\ObservedWollastonDenominator}$ Gran. Um die Dicke so äusserst zarter Gegenstände mit einem gewissen Grade von Genauigkeit messen zu können, bediente er sich einer besoneern Vorrichtung, welche er in dem andern Aufsatze unter dem Namen Single-lens micrometer beschreibt. Das Wesentliche besteht in einer Linse von sehr kurzer Brennweite ($\frac{1}{12}$ Zoll), in deren Brennpunkt das Object gebracht wird. Die Oeffnung dieser Linse ist so klein, dass neben der Linse in den Messingplättchen, in welches sie gefasst ist, in einer Entfernung von $\frac{1}{23}$ Zoll vom Mittelpunkt ein kleines Loch angebracht werden kann, und die Pupille zugleich von dem Object durch die Linse und von einer verschiebbaren Scale durch das Loch Licht erhält. Die Scale wird so weit entfernt, bis das Object genau einen oder einige Theile derselben deckt; die Entfernung der Scale und die Anzahl der Theile bestimmen dann die Grösse des Objects.

Noch gehört hieher: \emph{Ueber die Lichtstärke des Cassegrainschen Teleskop verglichen mit dem Gregoryschen} von H. \textsc{Kater}. Zwei Paare solcher Teleskope, von Einem Künstler (Hrn. \textsc{Crickmore} in Ipswich) verfertigt, wurden mit einander verglichen, und dabei die Verschiedenheit des Flächeninhalts der grossen Spiegel und der Abgang durch die kleinern Spiegel nebst den Armen, die sie trugen, in Rechnung gebracht. Bei dem einen Paar war das Verhältniss der Lichtstärke wie 7 zu 3, bei dem andern wie 3 zu 2, beide Male zu Gunsten des \textsc{Cassegrain}'schen Teleskops; beim ersten Paar hatte der Spiegel des \textsc{Cassegrain}'schen, beim andern der des \textsc{Gregory}'schen den Vorzug einer etwas vollkommnern Politur, daher der Verfasser glaubt, dass unter ganz gleichen Umständen das \textsc{Cassegrain}'sche Teleskop etwa doppelt so viel Licht habe, als das \textsc{Gregory}sche. Wenn wir voraussetzen dürfen, dass die Oeffnung, an welche unmittelbar das Auge angebracht wurde, in allen vier Teleskopen die gehörige Grösse hatte, um alles Licht von den grossen Spiegeln durchzulassen, — der Verfasser berührt diesen Punkt nicht, allein Rec. \SourceLeaf{568}hat verschiedene \textsc{Gregory}'sche Teleskope unter Händen gehabt, bei welchen diese Oeffnung zu klein war — so sind diese Erfahrungen sowohl dem Astronomen als dem Physiker höchst merkwürdig, und man möchte geneigt sein, der Erklärung des Verf. beizutreten, der den Grund des sonderbaren Phänomens in dem Umstande sucht, dass im \textsc{Gregory}'schen Teleskop zwei wirkliche Bilder zu Stande gebracht werden, im \textsc{Cassegrain}'schen hingegen nur eines, und daraus den Schluss zieht, dass das Kreuzen der Strahlen in einem Punkt oder die Formation eines wirklichen Bildes einen gewissen Lichtverlust bewirkt. Er glaubt, dass hiermit vielleicht die Erfahrung in Verbindung stehe, dass \textsc{Galilei}'sche Fernröhre unter gleichen Umständen mehr Wirkung thun, als astronomische.
